\section{Conclusion}
\label{sec:conclusion}

This paper develops a dynamic portfolio framework in which a fully neural, all-tree Student-$t$ D-vine and an LSTM-TD3 controller have distinct but connected roles. The D-vine generates causally ordered synthetic episodes, provides explicit conditional-dependence features, and supplies the scenarios used for CVaR. The agent selects self-financing long--short portfolios under gross and position limits and maximises dense increments of terminal-wealth CRRA utility after transaction and financing costs. The training protocol uses all synthetic data for pre-training, all permissible realised historical trajectories for one fine-tuning pass, and reserves the final 24 decisions for locked evaluation.

The frozen main result is competitive but mixed. Over 22 complete holding periods, the preregistered 20-seed mean-weight ensemble records the lowest observed volatility and highest observed Sharpe ratio among six benchmarks, while ranking third in terminal wealth and CRRA certainty-equivalent return. Its positive utility effect relative to equal weight is not statistically significant, all benchmark-effect intervals include zero, and individual-seed results are dispersed. Weight averaging removes most seed-specific short and leveraged exposure, so part of the ensemble's advantage is diversification and implementation-cost compression rather than a uniformly stable learned policy.

The post-holdout ablations impose an important limit on the novelty claim. A coarse no-vine TD3 ensemble performs better descriptively, and the pretrained-only ensemble is nearly indistinguishable from the fully fine-tuned policy on the consumed path. These results do not prove that vine information or fine-tuning is harmful, because the coarse intervention removes multiple channels and the sample is short. They do show that incremental component value has not yet been established. The preregistered matched-seed causal analysis therefore separates direct vine features, scenario-CVaR observation, CVaR reward shaping, synthetic pre-training, the neural generator, recurrence, fine-tuning, and the TD3 update rule.

Accordingly, the defensible contribution at this stage is methodological and empirical transparency: a dynamic dependence and control system, a validated synthetic-data protocol, publication-grade leakage and constraint controls, and a frozen mixed result that is not converted into a superiority narrative. Confirmation requires evidence independent of the consumed holdout. The next tests should use non-overlapping frozen walk-forward windows, an external investable total-return panel, and ultimately future observations collected after preregistration. Those designs must retain drift-aware turnover, actual-day financing, fail-closed optimisation, paired time-series inference, and complete source, data, environment, and retry provenance.
